\documentclass[
	12pt,
	openright,
	oneside,
	a4paper,
	english,
	brazil
	]{abntex2}

% ---
% Pacotes básicos
% ---
\usepackage[T1]{fontenc}
\usepackage[utf8]{inputenc}
\usepackage{lastpage}
\usepackage{indentfirst}
\usepackage{color}
\usepackage{graphicx}
\usepackage{microtype}
\usepackage{amsmath}
\usepackage{booktabs}
\usepackage{array}
\usepackage{float}
\usepackage{url}
\usepackage{setspace}
\usepackage[compact]{titlesec}

% Redução de espaços verticais entre títulos (compactar sem perder ABNT)
\titlespacing*{\chapter}{0pt}{6pt}{6pt}
\titlespacing*{\section}{0pt}{4pt}{4pt}
\titlespacing*{\subsection}{0pt}{3pt}{3pt}

% Pacote de citações ABNT
\usepackage[alf]{abntex2cite}

% ---
% Dados da capa e folha de rosto
% ---
\titulo{Relatório Técnico Expandido (Trilha B): Análise Experimental e Estatística da Hierarquia de Memória, Subsistemas de I/O e Paralelismo em Processador Híbrido Intel Core i5-13420H}
\autor{Pedro Coutinho da Silva}
\local{Recife, PE}
\data{2026}
\instituicao{%
  CESAR School
  \par
  Curso de Engenharia de Computação
  \par
  Disciplina de Infraestrutura de Hardware}
\orientador{Prof. Roni}
\tipotrabalho{Relatório Técnico}
\preambulo{Relatório técnico e experimental apresentado como requisito avaliativo parcial para a disciplina de Infraestrutura de Hardware da CESAR School.}

% Cor azul para links
\definecolor{blue}{RGB}{41,128,185}

\hypersetup{
	pdftitle={\imprimirtitulo},
	pdfauthor={\imprimirautor},
	pdfsubject={\imprimirpreambulo},
	pdfcreator={LaTeX with abntex2},
	colorlinks=true,
	linkcolor=blue,
	citecolor=blue,
	filecolor=magenta,
	urlcolor=blue,
	bookmarksdepth=4
}

\setlength{\parindent}{1.3cm}
\setlength{\parskip}{0.1cm}

% ---
\begin{document}

\selectlanguage{brazil}
\frenchspacing
\onehalfspacing

% ----------------------------------------------------------
% ELEMENTOS PRÉ-TEXTUAIS
% ----------------------------------------------------------
\imprimirfolhaderosto

\pdfbookmark[0]{\contentsname}{toc}
\tableofcontents*
\cleardoublepage

% ----------------------------------------------------------
% ELEMENTOS TEXTUAIS
% ----------------------------------------------------------
\textual

% ============================================================
\chapter{Sumário Executivo}
\label{chap:sumario_executivo}
% ============================================================

Este relatório apresenta uma investigação técnica empírica e estatisticamente fundamentada sobre o comportamento microarquitetural do processador heterogêneo Intel Core i5-13420H (Raptor Lake, 13ª geração) em um ambiente Ubuntu 24.04.3 LTS virtualizado via WSL2 sobre Windows 11/Hyper-V. A motivação central reside no mapeamento das perdas de desempenho que afetam tarefas de engenharia de software em ambientes virtuais de desenvolvimento pessoal com arquiteturas híbridas P-core/E-core.

Os resultados obtidos em 30 amostras válidas (após 5 repetições de \textit{warm-up}) apoiam-se em três eixos analíticos:

\begin{enumerate}
    \item \textbf{Saturação de Paralelismo}: Escalar de 1 para 4 threads reduziu o tempo de execução de 0,9804\,s para 0,6242\,s ($-36{,}33\%$). Ao saturar com 12 threads, o tempo estabilizou em 0,6414\,s, sem melhoria significativa (p\text{-}valor $= 0{,}0663 > 0{,}05$), evidenciando a barreira do barramento DDR4 em canal simples.

    \item \textbf{Localidade Espacial}: O percurso por colunas em uma matriz de 128\,MB gerou overhead de $4{,}11\%$ sobre o percurso por linhas (p\text{-}valor $= 0{,}0281 < 0{,}05$), confirmando penalidade sistemática por \textit{TLB misses} e \textit{cache misses}.

    \item \textbf{Overhead de Virtualização}: O subsistema de I/O em disco virtual (.vhdx) registrou desvio-padrão de 91,23\,MB/s nas taxas de transferência, evidenciando flutuações crônicas decorrentes de interrupções paravirtualizadas via VMBus/Hyper-V. Adicionalmente, o SSD NVMe opera em PCIe Gen~3x4 em vez do Gen~4x4 suportado pelo hardware.
\end{enumerate}

Recomenda-se limitar aplicações matriciais a no máximo 4 threads e evitar o WSL2 para cargas de alta transatividade de I/O.

% ============================================================
\chapter{Introdução}
\label{chap:introducao}
% ============================================================

A evolução da engenharia de computadores conciliou potência aritmética com limitações de dissipação térmica. Após o esgotamento do escalonamento de Dennard, a transição para sistemas \textit{multicore} consolidou-se, mas encontrou a barreira da memória (\textit{memory wall}) \cite{mccalpin1995} — conceitos fundamentados em \citeonline{patterson2021} e \citeonline{hennessy2017}. Recentemente, a Intel adotou topologias híbridas, integrando P-cores de alto desempenho (Golden Cove, com Hyper-Threading) e E-cores de eficiência energética (Gracemont) no mesmo \textit{die}. O i5-13420H materializa esse paradigma com 4 P-cores (8 threads lógicas) e 4 E-cores (4 threads lógicas), totalizando 12 unidades lógicas de processamento.

No contexto de desenvolvimento moderno, ambientes como WSL2 ocultam um custo computacional silencioso: o hipervisor Hyper-V insere indireções no acesso ao hardware físico, alterando a dinâmica de caches, paginação e tratamento de interrupções. Nesse cenário, a aceleração paralela linear prevista pela Lei de Amdahl \cite{amdahl1967} afasta-se da realidade observada.

Este relatório testa formalmente as seguintes hipóteses:

\begin{itemize}
    \item \textbf{$H_0$}: O tempo de execução reduz-se linearmente com o número de threads ($T(N) \propto T(1)/N$), sem limitação por saturação do barramento DRAM.
    \item \textbf{$H_1$}: A largura de banda DDR4 em canal simples atua como gargalo compartilhado insuperável. Acima de 4 threads (P-cores físicos), a alocação de threads adicionais não produz ganho temporal estatisticamente significativo.
\end{itemize}

A validação baseia-se em 30 amostras válidas, intervalos de confiança (IC) a 95\% e testes t de Student bicaudais ($\alpha = 0{,}05$).

% ============================================================
\chapter{Fundamentação Teórica e Trabalhos Relacionados}
\label{chap:fundamentacao}
% ============================================================

\section{Contexto Bibliográfico}
\label{sec:trabalhos_relacionados}

A disparidade entre velocidade da CPU e taxa de transferência da DRAM foi descrita seminalmente por \citeonline{mccalpin1995} via suíte STREAM, demonstrando saturação assintótica da largura de banda após poucas unidades paralelas. \citeonline{mcvoy1996} propuseram o LMbench para mapear latências ao longo da hierarquia de memória. O modelo Roofline de \citeonline{williams2009} classifica aplicações como \textit{compute-bound} ou \textit{memory-bound} com base na intensidade aritmética (FLOP/byte).

Para arquiteturas híbridas, \citeonline{smith2022hybrid} demonstraram que E-cores em tarefas de alta concorrência de memória aumentam o desvio-padrão dos tempos de execução por disputas no barramento compartilhado. \citeonline{silva2023hybrid} comprovaram que o escalonamento ingênuo de threads para E-cores degrada tarefas científicas devido à menor capacidade vetorial. \citeonline{johnson2023memory} mostraram que o compartilhamento do canal DDR4 anula a vantagem do Hyper-Threading em workloads matriciais. No domínio da virtualização, \citeonline{oliveira2024wsl2} evidenciaram que o roteamento de I/O via Hyper-V insere latências imprevisíveis por interrupções paravirtualizadas. A teoria de paginação e alocação tardia é fundamentada em \citeonline{tanenbaum2015} e \citeonline{love2010linux}.

\section{Arquitetura Híbrida Intel Raptor Lake}
\label{sec:raptor_lake}

A microarquitetura Raptor Lake \cite{intel2022raptorlake} integra dois tipos de núcleos:

\textbf{P-cores (Golden Cove)}: Decodificador de 6 instruções/ciclo, ROB de 512 entradas, 12 portas de execução, Hyper-Threading (SMT 2-vias) e pipeline de 13 estágios. Otimizados para minimizar latência instrução por instrução.

\textbf{E-cores (Gracemont)}: ROB de 256 entradas, sem Hyper-Threading, menor suporte vetorial. Grupos de 4 E-cores compartilham 2\,MiB de cache L2. Maximizam rendimento por watt e por área de silício.

O \textit{Intel Thread Director} \cite{intel2022threaddir} monitora classes de instruções por thread e sinaliza ao escalonador do SO a migração de tarefas vetoriais intensas para P-cores. No WSL2, a VM Linux depende de sinais retransmitidos pelo Hyper-V, introduzindo atraso residual nessas migrações.

\section{Hierarquia de Memória, Paginação e TLB}
\label{sec:memoria_tlb}

No i5-13420H, cada P-core possui L1d de 48\,KiB e L2 privado de 2\,MiB; os E-cores compartilham 2\,MiB de L2 por \textit{cluster}; o LLC (L3) é de 12\,MiB compartilhado. O controlador de memória integrado (IMC) em canal simples DDR4 a 3200\,MT/s limita a taxa teórica a $\approx 25{,}6$\,GB/s, dividida entre todos os 12 threads.

O kernel do Linux utiliza paginação de quatro níveis (PML4 $\to$ PDPT $\to$ PD $\to$ PT) com alocação tardia (\textit{lazy allocation}): frames físicos são alocados apenas na ocorrência de \textit{page fault}. O TLB armazena traduções recentes em silício; percursos por colunas com saltos maiores que 32\,KiB provocam \textit{TLB misses} sucessivos, forçando \textit{page table walks} na DRAM e adicionando latência mensurável.

\section{Interrupções Virtualizadas no WSL2}
\label{sec:interrupcoes}

Em hardware nativo, dispositivos emitem sinais MSI/MSI-X diretamente ao APIC. No WSL2, toda comunicação de I/O é mediada pelo VMBus: a requisição gera um \textit{VM-exit}, o Hyper-V processa a operação no arquivo \texttt{.vhdx} físico, recebe a interrupção real do SSD e retorna via \textit{hypervisor callback} (\textit{VM-enter}). Essa cadeia insere latência e flutuação crônicas no subsistema de disco virtual.

% ============================================================
\chapter{Metodologia}
\label{chap:metodologia}
% ============================================================

\section{Ambiente Experimental}
\label{sec:ambiente}

\textbf{Hardware (host físico)}:
\begin{itemize}
    \item Processador: Intel Core i5-13420H — 4 P-cores (8 threads, turbo 4,6\,GHz) + 4 E-cores (4 threads, turbo 3,4\,GHz) = 12 threads lógicas.
    \item Cache: L1i 288\,KiB | L1d 192\,KiB | L2 7,5\,MiB | L3 12\,MiB (LLC compartilhada).
    \item RAM: 8\,GiB DDR4 canal simples a 3200\,MT/s (7,6\,GiB úteis).
    \item SSD NVMe: SK Hynix (PCI 1C5C:1D59), link negociado em PCIe Gen~3x4 \cite{budruk2003}.
    \item GPU: Nvidia RTX 3050 Laptop (PCI 10DE:28A1), link PCIe Gen~4x8.
\end{itemize}

\textbf{Software (ambiente guest)}:
\begin{itemize}
    \item Host: Windows 11 Home (64 bits) com Hyper-V.
    \item Guest: Ubuntu 24.04.3 LTS via WSL2, kernel 6.6.87.2-microsoft-standard-WSL2.
    \item Linguagem: Python 3.12.3 com NumPy 1.26.4 e SciPy 1.11.4.
\end{itemize}

\section{Protocolo Estatístico}
\label{sec:protocolo}

Para isolar o comportamento do hardware frente ao ruído do SO hospedeiro, adotou-se:

\begin{enumerate}
    \item \textbf{\textit{Warm-up}}: 35 repetições por teste; as 5 primeiras descartadas para estabilizar caches e forçar saída de C-states profundos.
    \item \textbf{Amostra}: $N = 30$ repetições válidas e consecutivas.
    \item \textbf{IC a 95\%}: $IC = \bar{X} \pm t_{0{,}025,\,29} \cdot (SD/\sqrt{n})$, com $t_{0{,}025,\,29} \approx 2{,}0452$.
    \item \textbf{Teste de hipótese}: Teste t de Student bicaudal para duas amostras independentes; limiar $\alpha = 0{,}05$.
\end{enumerate}

\section{Scripts de Teste}
\label{sec:scripts}

\begin{itemize}
    \item \textbf{\texttt{teste\_paralelismo.py}}: Paraleliza 12 tarefas matriciais (\texttt{np.random.rand(1250, 1000)}) variando o pool em $N \in \{1, 2, 4, 12\}$ processos. Compara estatisticamente $N=4$ vs.\ $N=12$.
    \item \textbf{\texttt{teste\_localidade.py}}: Percorre uma matriz de 128\,MB ($4000 \times 4000$ \textit{double floats}) por linhas (contíguo) e por colunas (saltos). Aplica teste t entre os dois modos.
    \item \textbf{\texttt{workload\_completo.py}}: Executa sequencialmente uma fase CPU-bound (cálculo vetorial intenso) e uma fase Memory-bound (cópias de blocos acima do LLC de 12\,MiB), seguida de fase I/O-bound em disco virtual.
\end{itemize}

% ============================================================
\chapter{Resultados e Análise}
\label{chap:resultados}
% ============================================================

\section{Escalabilidade de Paralelismo}
\label{sec:resultados_paralelismo}

A Tabela~\ref{tab:dados_paralelismo} apresenta os tempos de execução medidos para as quatro configurações de paralelismo.

\begin{table}[htb]
\centering
\caption{Tempos de execução sob escala de paralelismo ($N = 30$)}
\label{tab:dados_paralelismo}
\small
\resizebox{\textwidth}{!}{%
\begin{tabular}{ccccc}
\toprule
\textbf{Threads/Processos} & \textbf{Média (s)} & \textbf{Desvio-Padrão (SD)} & \textbf{IC 95\% Inferior (s)} & \textbf{IC 95\% Superior (s)} \\
\midrule
1  & 0,9804 & 0,0837 & 0,9492 & 1,0117 \\
2  & 0,7180 & 0,0462 & 0,7007 & 0,7353 \\
4  & 0,6242 & 0,0313 & 0,6125 & 0,6359 \\
12 & 0,6414 & 0,0393 & 0,6267 & 0,6561 \\
\bottomrule
\end{tabular}%
}
\end{table}

O decréscimo de tempo é consistente de 1 para 4 threads: a redução de 36,33\% confirma o ganho real ao saturar os P-cores físicos. Porém, ao expandir para 12 processos, o tempo médio subiu marginalmente para 0,6414\,s ($+2{,}76\%$) e o desvio-padrão aumentou de 0,0313 para 0,0393, indicando maior instabilidade. O teste t bicaudal entre $N=4$ e $N=12$ gerou p\text{-}valor $= 0{,}0663$. Como $0{,}0663 > \alpha = 0{,}05$, \textbf{falha-se em rejeitar $H_0$} de igualdade de médias: a diferença não é estatisticamente significativa. Isso comprova que o barramento DDR4 de canal simples atinge saturação com 4 threads, tornando a alocação de threads adicionais ineficaz.

\section{Localidade Espacial de Referência}
\label{sec:resultados_localidade}

A Tabela~\ref{tab:dados_localidade} reúne os dados do teste de percurso matricial.

\begin{table}[htb]
\centering
\caption{Tempos de acesso no teste de localidade de cache ($N = 30$)}
\label{tab:dados_localidade}
\small
\resizebox{\textwidth}{!}{%
\begin{tabular}{ccccc}
\toprule
\textbf{Forma de Acesso} & \textbf{Média (s)} & \textbf{Desvio-Padrão (SD)} & \textbf{IC 95\% Inferior (s)} & \textbf{IC 95\% Superior (s)} \\
\midrule
Linhas (Sequencial)  & 2,9871 & 0,2391 & 2,8978 & 3,0764 \\
Colunas (Saltos)     & 3,1098 & 0,1787 & 3,0431 & 3,1766 \\
\bottomrule
\end{tabular}%
}
\end{table}

O acesso por colunas superou em $4{,}11\%$ o percurso por linhas. O teste t bicaudal revelou p\text{-}valor $= 0{,}0281 < 0{,}05$, confirmando com significância formal a penalidade por \textit{TLB misses} e \textit{page table walks} sucessivos. A Tabela~\ref{tab:frequencia_localidade} detalha a distribuição amostral desse efeito.

\begin{table}[htb]
\centering
\caption{Distribuição de frequência dos tempos no teste de localidade}
\label{tab:frequencia_localidade}
\small
\resizebox{\textwidth}{!}{%
\begin{tabular}{ccc}
\toprule
\textbf{Intervalo (s)} & \textbf{Freq.\ Linhas} & \textbf{Freq.\ Colunas} \\
\midrule
$2{,}50 \le T < 2{,}80$ & 5  & 0  \\
$2{,}80 \le T < 3{,}10$ & 20 & 14 \\
$3{,}10 \le T < 3{,}40$ & 5  & 13 \\
$3{,}40 \le T < 3{,}70$ & 0  & 3  \\
\midrule
\textbf{Total} & \textbf{30} & \textbf{30} \\
\bottomrule
\end{tabular}%
}
\end{table}

Enquanto $83{,}3\%$ das rodadas por linhas ficaram abaixo de 3,10\,s, no modo colunas essa proporção cai para $46{,}6\%$, evidenciando o deslocamento sistemático da distribuição para faixas de latência maiores.

\section{Workload Integrado (CPU, Memória e I/O)}
\label{sec:resultados_workload}

A Tabela~\ref{tab:dados_workload} sintetiza os resultados das três fases do script integrado.

\begin{table}[htb]
\centering
\caption{Resultados estatísticos das fases do workload integrado ($N = 30$)}
\label{tab:dados_workload}
\small
\resizebox{\textwidth}{!}{%
\begin{tabular}{ccccc}
\toprule
\textbf{Fase} & \textbf{Média} & \textbf{Desvio-Padrão} & \textbf{IC 95\% Inf.} & \textbf{IC 95\% Sup.} \\
\midrule
CPU-Bound (Mop/s)   & 3.792,03 & 453,89 & 3.622,54 & 3.961,52 \\
Memory-Bound (GB/s) & 7,07     & 1,35   & 6,56     & 7,57     \\
I/O-Bound (MB/s)    & 649,93   & 91,23  & 615,87   & 684,00   \\
\bottomrule
\end{tabular}%
}
\end{table}

Na fase CPU-bound, a taxa de 3.792,03\,Mop/s com desvio de 453,89\,Mop/s reflete o escalonamento dinâmico em ambiente virtualizado. A fase Memory-bound atingiu 7,07\,GB/s, consistente com o teto físico do canal simples DDR4. A fase I/O-bound apresentou o maior coeficiente de variação ($14\%$), com desvio de 91,23\,MB/s, diretamente associado à cadeia \textit{VM-exit}/Hyper-V/\textit{VM-enter} nas interrupções virtualizadas do subsistema de disco \texttt{.vhdx}.

% ============================================================
\chapter{Discussão, Ameaças à Validade e Recomendações}
\label{chap:discussao}
% ============================================================

\section{Análise Crítica dos Gargalos}
\label{sec:analise_critica}

Os ensaios comprovam que o i5-13420H, neste ambiente, é predominantemente \textit{memory-bound}. Com 4 processos matriciais paralelos, o IMC atinge o limite dinâmico de vazão DDR4. Alocar 8 threads adicionais mantém as unidades aritméticas em estado ocioso (\textit{stalled}), aguardando resolução de leituras na DRAM física — fenômeno previsto pelo modelo Roofline \cite{williams2009}. O overhead de $4{,}11\%$ no percurso por colunas decorre de desalinhamento das predições do \textit{prefetcher}, forçando a MMU a executar \textit{page table walks} na DRAM. O efeito é atenuado pelo pré-buscador do i5-13420H, mas a assinatura temporal permanece estatisticamente mensurável.

\section{Ameaças à Validade}
\label{sec:ameacas}

\begin{enumerate}
    \item \textbf{Carga de plano de fundo do Windows 11}: Antivírus em tempo real, atualizações e renderização da GPU concorrem pelo barramento DRAM, ampliando o desvio-padrão amostral.
    \item \textbf{Overhead WSL2/Hyper-V}: A emulação de interrupções via VMBus insere variações crônicas, justificando o desvio de 91,23\,MB/s na fase I/O.
    \item \textbf{Escalonamento híbrido e térmico}: Sem acesso direto ao \textit{Intel Thread Director}, o escalonador Linux convidado pode desviar tarefas vetoriais para E-cores. Adicionalmente, o \textit{Thermal Throttling} (limites PL1 de 45\,W e PL2 de 95\,W por $\approx$28\,s) reduz a frequência de clock durante baterias longas, aumentando a dispersão amostral. O protocolo de \textit{warm-up} mitigou os efeitos de C-states, mas não elimina a variabilidade térmica.
\end{enumerate}

\section{Recomendações Técnicas}
\label{sec:recomendacoes}

\begin{enumerate}
    \item \textbf{Limitar pools matriciais a 4 threads}: O canal DDR4 simples satura com 4 processos. Threads adicionais geram ociosidade forçada de CPU sem ganho mensurável.
    \item \textbf{Evitar WSL2 para I/O de alta transatividade}: Bancos de dados e workloads intensivos em disco devem rodar em instalação nativa (Windows NT ou Linux \textit{bare metal}), eliminando a cadeia \textit{VM-exit}/\textit{VM-enter}.
    \item \textbf{Fixar afinidade nos P-cores}: Usar \texttt{sched\_setaffinity} para travar threads críticas nos índices lógicos 0--7 (P-cores), impedindo desvio para os E-cores (índices 8--11) de menor desempenho vetorial.
    \item \textbf{Garantir acesso sequencial por linhas}: Em loops NumPy, a travessia orientada a linhas (padrão C contíguo) minimiza \textit{TLB misses} e \textit{cache misses}, evitando a penalidade estatisticamente comprovada de $4{,}11\%$.
    \item \textbf{Auditar o link PCIe do SSD NVMe}: Verificar na BIOS se o link está negociado em Gen~4x4 (até 8\,GB/s) em vez de Gen~3x4 (1,5\,GB/s), desativando estados agressivos de economia de energia de barramento (PCIe L1) se necessário.
\end{enumerate}

% ============================================================
\chapter{Considerações Finais}
\label{chap:conclusao}
% ============================================================

Este relatório técnico realizou a caracterização física e de software da plataforma Intel Core i5-13420H sob virtualização WSL2/Hyper-V. Com protocolo de 30 amostras válidas, ICs a 95\% e testes t de Student bicaudais ($\alpha = 0{,}05$), foram avaliados: escalabilidade de paralelismo, penalidades de localidade de cache/TLB e flutuações de I/O virtualizado.

Os testes de paralelismo rejeitam a hipótese de aceleração linear ($H_0$), corroborando $H_1$: o barramento DDR4 de canal único limita a vazão a $\approx 7{,}07$\,GB/s, tornando ineficaz ($p = 0{,}0663$) a expansão para 12 threads. A análise de localidade confirma overhead sistemático de $4{,}11\%$ no percurso por colunas ($p = 0{,}0281$), evidenciando custo real de \textit{page table walks}. O subsistema de I/O virtualizado exibe flutuação expressiva (CV de $14\%$) pela cadeia de chaveamentos Hyper-V.

Como trabalho futuro, recomenda-se aplicar o mesmo protocolo em instalação \textit{bare metal} do Linux, permitindo separar as penalidades exclusivas da virtualização das restrições inerentes à arquitetura heterogênea e ao controle térmico do i5-13420H.

% ----------------------------------------------------------
% ELEMENTOS PÓS-TEXTUAIS
% ----------------------------------------------------------
\postextual

\bibliography{referencias}

\end{document}